\section{Conclusion and discussion}

Risks in LLMs involve possible errors in answering questions with high confidence.
One possible way to mitigate this problem is by highlighting those questions that, with high probability, would be answered incorrectly.
We suggest using uncertainty estimations for LLMs to identify possible errors.
For a more detailed analysis, the paper examines various subject areas, question types, and specific requirements for knowledge and reasoning.
We explored how different uncertainty estimates for LLMs highlight potential errors of a model and examined various aspects of problems in QA to identify what type of problems a specific uncertainty estimate can highlight.
Our focus was on reliable and popular methods, including a model as a judge and an entropy-based approach.

Experimental evidence suggests that different uncertainty estimates perform differently across diverse subject areas.
Our findings also demonstrate that the key factor is the varying composition of question sets within a subject area collected for a specific subject area: entropy performs much better when no reasoning or little reasoning is required. Thus, it primarily highlights the lack of the required knowledge. 
The conclusion also depends on the model used: for reasoning-oriented, such as Qwen-72B, entropy better highlights the possibility of an error compared to Mistral.